\begin{abstract}
Test-time adaptation (TTA) updates deployed models from unlabeled target inputs, but evidence that is useful in one deployment setting can become unreliable under a different stream, model interface, or adaptation entity. We present \textbf{A}daptive \textbf{T}est-time \textbf{L}earning \textbf{A}cross \textbf{S}cenarios (\textbf{ATLAS}), a reliability-oriented TTA formulation organized around three design principles. \textbf{CORE} selects image, pixel, or prompt-view entities only when source, raw, and semantics-preserving views provide corroborating evidence while a destroy view tests structural sensitivity. \textbf{CARE} anchors accepted updates using either source-relative clipping or class-centered support. \textbf{SANE} normalizes the loss by selected entities rather than the raw tensor size. We instantiate these principles with task-specific parameter scopes, gates, and optimization protocols for corruption robustness, natural shifts, non-i.i.d. streams, dense prediction, and episodic prompt adaptation. Under this five-level evaluation, ATLAS reaches $68.8\%$ on CIFAR-100-C continual adaptation, $65.14\%$ on ViT-B/16 ImageNet-C continual adaptation, and $64.2\%$ OOD average accuracy for CLIP ViT-B/16 prompt adaptation with CoOp initialization. Component ablations attribute the classification gains to evidence selection and anchoring, while the segmentation ablation shows that selected-entity normalization is necessary in the evaluated dense-prediction setting. These results support the three reliability principles across the studied interfaces without assuming one numerical recipe for every scenario. Part of the core code is included in the supplementary materials.
\end{abstract}
